% GENERATED FILE. Edit canonical lessons, then run npm run generate:books.
% canonical-source-hash: fnv1a64:e81f597c74c71154
% canonical-lessons: KA-C60-cow, KA-C60-goat, KA-C60-hen, KA-C60-crow, KA-C60-bird

\chapter{Animals of the House}
\label{ch:ka-animals}
\hlchaptermodality{60}

\begin{chapteropening}
\textbf{By the end of this chapter:} \emph{I can name five animals a Kannada household lives with, from the cow to the bird.}

Complete KA-C60-bird and say all five words in order.
\end{chapteropening}

\section[hasu]{\kn{ಹಸು} (hasu) — a cow}

\label{lesson:KA-C60-cow}



\begin{quote}
Before the new run: what did \emph{kesaru} mean, and what did \emph{beṅki} mean?
\end{quote}

\subsection*{You'll want to know: \kn{ಹಸು}}
\textbf{\kn{ಹಸು}} (\emph{hasu}) — \textquotedblleft{}a cow\textquotedblright{}.

\kn{ಹಸು} is inherited, and it is the animal a Karnataka household is built around. The \kn{ಹಾಲು} you met in the drinks chapter comes out of it, and most of what a Kannada kitchen makes begins with what comes out of it.

It belongs to the \kn{ರೈತ} before it belongs to anybody, and in a \kn{ಹಳ್ಳಿ} it is tied in the yard rather than kept in a shed, close enough to the \kn{ಬಾಗಿಲು} that you step around it going in.

\begin{grammarlens}[title={What you've built}]
The first of five animals.
\end{grammarlens}

\subsection*{Your turn}
Say these aloud:

\begin{itemize}
\raggedright
  \item \emph{hasu}
  \item it once more, slowly
  \item \emph{hasu}, then \emph{hālu}, and say which one comes out of the other
\end{itemize}

\begin{tcolorbox}[breakable,colback=teal!4,colframe=teal!35!black,title={Before you move on}]
What does \kn{ಹಸು} mean? (\textquotedblleft{}a cow\textquotedblright{}.) Say it once more.
\end{tcolorbox}
\section[mēke]{\kn{ಮೇಕೆ} (mēke) — a goat}

\label{lesson:KA-C60-goat}



\begin{quote}
Before the new one: what did \emph{hasu} mean?
\end{quote}

\subsection*{You'll want to know: \kn{ಮೇಕೆ}}
\textbf{\kn{ಮೇಕೆ}} (\emph{mēke}) — \textquotedblleft{}a goat\textquotedblright{}.

\kn{ಮೇಕೆ} is shared right across the family — Telugu says it almost exactly the same way — and it names the animal that goes where the \kn{ಹಸು} does not.

A \kn{ಹಸು} is walked to flat ground. A \kn{ಮೇಕೆ} is let up a \kn{ಬೆಟ್ಟ}, and it will strip the \kn{ಎಲೆ} off a low \kn{ಕೊಂಬೆ} standing on its back legs to reach. That difference is the whole reason a household keeps both.

\begin{grammarlens}[title={What you've built}]
Two.
\end{grammarlens}

\subsection*{Your turn}
Say these aloud:

\begin{itemize}
\raggedright
  \item \emph{mēke}
  \item it once more, slowly
  \item \emph{hasu}, then \emph{mēke}, and say which one climbs
\end{itemize}

\begin{tcolorbox}[breakable,colback=teal!4,colframe=teal!35!black,title={Before you move on}]
What does \kn{ಮೇಕೆ} mean? (\textquotedblleft{}a goat\textquotedblright{}.) Say it once more.
\end{tcolorbox}
\section[kōḷi]{\kn{ಕೋಳಿ} (kōḷi) — a hen}

\label{lesson:KA-C60-hen}



\begin{quote}
Before the new one: what did \emph{mēke} mean?
\end{quote}

\subsection*{You'll want to know: \kn{ಕೋಳಿ}}
\textbf{\kn{ಕೋಳಿ}} (\emph{kōḷi}) — \textquotedblleft{}a hen\textquotedblright{}.

\kn{ಕೋಳಿ} is the domestic fowl, hen or cock alike; the word does not divide them, and where the difference matters a second word goes in front of it.

It is also the clock of the \kn{ಹಳ್ಳಿ}. A \kn{ಕೋಳಿ} wakes a house before \kn{ಬೆಳಗ್ಗೆ} is properly under way, which is a large part of why nobody there needed a \kn{ಗಂಟೆ} to do the same job.

\begin{grammarlens}[title={What you've built}]
Three.
\end{grammarlens}

\subsection*{Your turn}
Say these aloud:

\begin{itemize}
\raggedright
  \item \emph{kōḷi}
  \item it once more, slowly
  \item \emph{kōḷi}, then \emph{beḷagge}, and say which one comes first
\end{itemize}

\begin{tcolorbox}[breakable,colback=teal!4,colframe=teal!35!black,title={Before you move on}]
What does \kn{ಕೋಳಿ} mean? (\textquotedblleft{}a hen\textquotedblright{}.) Say it once more.
\end{tcolorbox}
\section[kāge]{\kn{ಕಾಗೆ} (kāge) — a crow}

\label{lesson:KA-C60-crow}



\begin{quote}
Before the new one: what did \emph{kōḷi} mean?
\end{quote}

\subsection*{You'll want to know: \kn{ಕಾಗೆ}}
\textbf{\kn{ಕಾಗೆ}} (\emph{kāge}) — \textquotedblleft{}a crow\textquotedblright{}.

\kn{ಕಾಗೆ} is an imitation word: the bird is named for the noise it makes, and the same trick produced the crow-word in each of the sister languages. Once you hear one you have the etymology.

Of all the birds in this chapter it is the one a household actually deals with day to day. Cooked rice set out on a step at the back of the house is put there for the \kn{ಕಾಗೆ}, and it arrives within the minute.

\begin{grammarlens}[title={What you've built}]
Four.
\end{grammarlens}

\subsection*{Your turn}
Say these aloud:

\begin{itemize}
\raggedright
  \item \emph{kāge}
  \item it once more, slowly
  \item \emph{kāge}, then \emph{kōḷi}, and say which one is kept and which one turns up
\end{itemize}

\begin{tcolorbox}[breakable,colback=teal!4,colframe=teal!35!black,title={Before you move on}]
What does \kn{ಕಾಗೆ} mean? (\textquotedblleft{}a crow\textquotedblright{}.) Say it once more.
\end{tcolorbox}
\section[hakki]{\kn{ಹಕ್ಕಿ} (hakki) — a bird}

\label{lesson:KA-C60-bird}



\begin{quote}
Before the new one: what did \emph{kāge} mean?
\end{quote}

\subsection*{You'll want to know: \kn{ಹಕ್ಕಿ}}
\textbf{\kn{ಹಕ್ಕಿ}} (\emph{hakki}) — \textquotedblleft{}a bird\textquotedblright{}.

\kn{ಹಕ್ಕಿ} is the general one: any bird at all, the \kn{ಕೋಳಿ} and the \kn{ಕಾಗೆ} both inside it. It is Kannada's own, and it sits beside a borrowed word that covers the same ground more formally — the arrangement you already saw when \kn{ಸ್ನೇಹಿತ} turned out to have \kn{ಗೆಳೆಯ} living quietly beside it.

One warning before you say it out loud. Written in our letters, \emph{hakki} is \emph{akki} with an \emph{h} on the front, and \kn{ಅಕ್ಕಿ} is the raw rice from the water-and-rice chapter. In Kannada letters they start differently and there is no danger; in a romanization there is. Say the \emph{h}.

\begin{grammarlens}[title={What you've built}]
Five: \kn{ಹಸು}, \kn{ಮೇಕೆ}, \kn{ಕೋಳಿ}, \kn{ಕಾಗೆ}, \kn{ಹಕ್ಕಿ}. Enough to greet a house and name what lives in it.
\end{grammarlens}

\subsection*{Your turn}
Say these aloud:

\begin{itemize}
\raggedright
  \item \emph{hakki}
  \item it once more, slowly
  \item all five in order, then \emph{namaskāra} to whoever keeps them
\end{itemize}

\begin{tcolorbox}[breakable,colback=teal!4,colframe=teal!35!black,title={Before you move on}]
What does \kn{ಹಕ್ಕಿ} mean? (\textquotedblleft{}a bird\textquotedblright{}.) Say it once more.
\end{tcolorbox}
